\chapter{Preliminaries}

\dropcap{I}{n this chapter} we review the basic language of algebraic topology and differential geometry that is needed to discuss Seiberg-Witten theory and the Donaldson theorem. In the differential realm, we assume familiarity with smooth manifolds, de Rham cohomology, vector bundles, and some Riemannian geometry. In algebraic topology, we assume familiarity with singular and cellular homology.

\section{Gauge theory}
Gauge theory is the heart of Seiberg-Witten theory. The name comes from physics, where gauge fields are fields with symmetries and redundant degrees of freedom. Mathematically, the study of gauge fields is done in terms of principal bundles, connections, curvature, and associated vector bundles.

\subsection{Principal bundles, connections and curvature}
The basic objects of gauge theory are principal bundles.

\begin{definition}[Principal $G$-bundle]
Let $M$ be a smooth manifold and let $G$ be a Lie group. A principal $G$-bundle over $M$ is a smooth manifold $P$ together with a smooth surjective map $\pi\colon P\to M$ and a right action of $G$ on $P$ satisfying equivariance and local triviality conditions.
\end{definition}

We write this compactly as
\[
  G \hookrightarrow P \xrightarrow{\pi} M.
\]

\begin{example}[Frame bundles]
Let $E\to M$ be a $\K$-vector bundle of rank $k$. The frame bundle $\operatorname{Fr}(E)$ is the bundle whose fiber over $x\in M$ consists of all ordered bases of $E_x$. It is naturally a principal $\operatorname{GL}(k,\K)$-bundle.
\[
  \operatorname{Fr}(E)=\coprod_{x\in M}\operatorname{Fr}_x(E).
\]
If $E$ is a real Riemannian vector bundle, one obtains the orthogonal frame bundle $O(E)$. If $E$ is a complex Hermitian bundle, one obtains the unitary frame bundle $U(E)$.
\end{example}

\begin{definition}[Connection on a principal bundle]
Let $G$ be a Lie group with Lie algebra $\g$. A connection on a principal bundle $G\hookrightarrow P\to M$ is a $\g$-valued one-form $\omega\in\Omega^1(P,\g)$ satisfying the standard equivariance and normalization conditions.
\end{definition}

\begin{proposition}[Transformation of local gauge potentials]
Let $\omega$ be a connection on a principal bundle $G\hookrightarrow P\to M$, and let $\{U_i\}_{i\in I}$ be a trivializing cover with local gauges. The local gauge potentials are related on overlaps by
\[
  \mathcal{A}_j = \Ad_{g_{ij}^{-1}}\circ \mathcal{A}_i + g_{ij}^{*}\Theta.
\]
If $G$ is a matrix Lie group, this becomes
\[
  \mathcal{A}_j = g_{ij}^{-1}\mathcal{A}_i g_{ij}+g_{ij}^{-1}\dd g_{ij}.
\]
\end{proposition}

\begin{definition}[Curvature]
The curvature of a connection $\omega$ is the two-form
\[
  \Omega := \dd^\omega \omega.
\]
Equivalently, it satisfies Cartan's structure equation
\[
  \Omega = \dd\omega + \frac12[\omega,\omega].
\]
\end{definition}
\subsection{Associated bundles and matter fields}
Let $G\hookrightarrow P\to M$ be a principal bundle, $V$ a vector space, and $\rho\colon G\to \operatorname{GL}(V)$ a representation. The associated bundle is
\[
  P\times_\rho V = P\times V/G.
\]
Sections of this bundle can be interpreted as matter fields in the language of physics.

\subsection{The space of connections}

\begin{proposition}[Connections form an affine space]
Let $\Conn(P)$ be the set of connections of a principal bundle $G\hookrightarrow P\to M$. Then $\Conn(P)$ is an affine space modeled on the vector space of Ad-tensorial one-forms $\Omega^1_{\Ad}(P,\g)$.
\end{proposition}

\section{Algebraic topology}
The main algebraic-topological tools needed later are cohomology, characteristic classes, orientations, and the intersection form of four-manifolds.

\subsection{Cohomology and cups}
\begin{definition}[Pairing of cohomology and homology]
The natural pairing of cohomology and homology is
\[
  \langle -,-\rangle\colon H^k(C,G)\otimes H_k(C)\to G,
\]
given by evaluating a cohomology class on a homology class.
\end{definition}

\begin{theorem}[Universal coefficients for cohomology]
Let $C_\bullet$ be a chain complex of free abelian groups and let $G$ be an abelian group. Then there is a split exact sequence
\[
0\to \Ext(H_{k-1}(C),G)\to H^k(C,G)\to \Hom(H_k(C),G)\to 0.
\]
\end{theorem}

\subsection{Orientations and caps}
A local orientation at a point of an $n$-manifold is a choice of generator of the local homology group. A global orientation is a locally consistent family of such choices.

\subsection{The intersection pairing}
The intersection pairing on a closed oriented four-manifold is one of the central objects in this story. It may be defined using Poincaré duality and the cup product.

\section{Characteristic classes}
Characteristic classes measure obstructions to trivializing bundles and will later be used to formulate the existence of spin and $\Spinc$ structures.

\subsection{Classifying spaces for line bundles and the first Chern class}
Complex line bundles are classified by maps into $\CP^\infty$, and the first Chern class records the corresponding cohomology class.

\subsection{Čech cohomology and the first Stiefel-Whitney class}
The first Stiefel-Whitney class detects the obstruction to orientability.

\section{Hodge theory}
This section would collect the Hodge-theoretic facts used later: harmonic representatives, Hodge decomposition, and self-dual forms.

\section{Elliptic regularity}
\subsection{Sobolev spaces}
Sobolev spaces provide the functional-analytic setting for elliptic differential operators.

\subsection{Elliptic operators}
Elliptic regularity is one of the analytic mechanisms behind smoothness statements for moduli spaces.
